\documentclass[11pt]{article}
% Shared preamble for Theseus user guides.
% Compiled with pdflatex. See docs/guides/BUILD.md for rationale.

\usepackage[a4paper,margin=0.85in]{geometry}
\usepackage[T1]{fontenc}
\usepackage[utf8]{inputenc}

% Body: Palatino (via mathpazo). Headings: Helvetica (sans).
% Palatino is requested in the house style as a substitute for Charter
% because Charter is not part of the standard TeX Live distribution
% that the build container ships.
\usepackage{mathpazo}
\usepackage[scaled=0.92]{helvet}
\usepackage{courier}
\renewcommand{\familydefault}{\rmdefault}

\usepackage{microtype}
\usepackage{parskip}
\usepackage{enumitem}
\usepackage{booktabs}
\usepackage{longtable}
\usepackage{titlesec}
\usepackage{xcolor}
\usepackage{fancyhdr}
\usepackage{fancyvrb}
\usepackage{graphicx}
% Allow `_` in text mode without escaping. The screenshot file names and
% some environment-variable references in the guides contain underscores.
\usepackage{underscore}
\usepackage{tcolorbox}
\tcbuselibrary{breakable,skins}
\usepackage[hidelinks,colorlinks=true,linkcolor=black,urlcolor=teal,citecolor=black]{hyperref}

% Sans for chapter and section headings; serif body keeps prose readable
% on screen and in print.
\titleformat{\section}{\sffamily\Large\bfseries}{\thesection}{0.7em}{}
\titleformat{\subsection}{\sffamily\large\bfseries}{\thesubsection}{0.6em}{}
\titleformat{\subsubsection}{\sffamily\normalsize\bfseries}{\thesubsubsection}{0.5em}{}

\pagestyle{fancy}
\fancyhf{}
\fancyhead[L]{\small\sffamily Theseus User Guide}
\fancyhead[R]{\small\sffamily\thepage}
\renewcommand{\headrulewidth}{0.3pt}

\setlength{\parindent}{0pt}

% Convenience commands shared by every guide.
\newcommand{\page}[1]{\textbf{#1}}
\newcommand{\file}[1]{\texttt{#1}}
\newcommand{\route}[1]{\texttt{#1}}
\newcommand{\env}[1]{\texttt{#1}}
\newcommand{\code}[1]{\texttt{#1}}

% Callout boxes used by every guide.
\newtcolorbox{tldr}{
  colback=teal!5!white,colframe=teal!55!black,
  fonttitle=\sffamily\bfseries,title=TL;DR,
  breakable,boxrule=0.4pt,arc=2pt
}
\newtcolorbox{youwillneed}{
  colback=gray!4!white,colframe=gray!55!black,
  fonttitle=\sffamily\bfseries,title=You will need,
  breakable,boxrule=0.4pt,arc=2pt
}
\newtcolorbox{notcovered}{
  colback=orange!5!white,colframe=orange!55!black,
  fonttitle=\sffamily\bfseries,title=This guide does NOT cover,
  breakable,boxrule=0.4pt,arc=2pt
}
\newtcolorbox{preview}{
  colback=yellow!8!white,colframe=yellow!55!black,
  fonttitle=\sffamily\bfseries,title=Preview --- partially shipped,
  breakable,boxrule=0.4pt,arc=2pt
}
\newtcolorbox{nextup}{
  colback=blue!4!white,colframe=blue!55!black,
  fonttitle=\sffamily\bfseries,title=Where to read next,
  breakable,boxrule=0.4pt,arc=2pt
}

% Insert a screenshot only when the file exists, so the build does not
% break before the Playwright capture run is wired up.
\newcommand{\screenshot}[3]{%
  \IfFileExists{screenshots/#1}{%
    \begin{figure}[h]\centering
      \includegraphics[width=\linewidth]{screenshots/#1}
      \caption{#2}\label{fig:#3}
    \end{figure}%
  }{%
    \begin{center}\small\itshape%
      [screenshot \texttt{screenshots/#1} not yet captured --- #2]%
    \end{center}%
  }%
}


\title{\sffamily\bfseries Theseus --- Quick Start\\[2pt]\large Guide 1 of 6}
\date{}
\author{}

\begin{document}
\maketitle
\thispagestyle{empty}

\begin{tldr}
Theseus is three programs that work together. \page{Dialectic} listens to
your conversations on the desktop and writes a structured transcript.
\page{Noosphere} is the workshop --- a background engine --- that turns
recorded material into a navigable library of claims, principles, and
conclusions. \page{Theseus Codex} is the website where you sign in,
upload material, review what the workshop has produced, and explore the
firm's working memory. By the end of this guide (about forty minutes)
you will have signed in, taken a short walk through the dashboard,
recorded a five-minute Dialectic session, watched the analysis land in
the Codex, uploaded one document, browsed the resulting claims and
conclusions, asked the Oracle a question, and seen how a Currents
opinion is formed.
\end{tldr}

\begin{youwillneed}
\begin{itemize}[leftmargin=*,itemsep=0.2em,topsep=0.2em]
  \item Your firm credentials. Your account was seeded for you; there is
        no public signup. If you do not have a working sign-in yet,
        contact whoever onboarded you.
  \item A browser pointed at the Codex (\route{www.theseuscodex.com} or
        your firm's deployment URL).
  \item The Dialectic desktop app installed (\route{Dialectic.dmg} for
        macOS or \route{Dialectic-Setup.exe} for Windows; both are on
        the firm's GitHub releases page).
  \item A working microphone.
  \item One short document you are willing to upload as the first
        corpus item --- a memo, transcript, essay, or PDF. Public
        material is easiest for a first walkthrough.
  \item Roughly forty uninterrupted minutes.
\end{itemize}
\end{youwillneed}

\begin{notcovered}
\begin{itemize}[leftmargin=*,itemsep=0.2em,topsep=0.2em]
  \item How extraction, clustering, principles, and conclusions work
        inside the workshop --- see Guide 2.
  \item The Oracle's citation contract, abstention, and hallucination
        warning in depth --- Guide 3.
  \item How Currents pulls in events, gates them, and generates
        opinions --- Guide 4.
  \item Prediction markets, the Market Decision Trace, paper vs.\ live,
        and the eight-gate safety architecture --- Guide 5.
  \item Operator-only surfaces (publication approval, methodology
        review week, kill switches, deletion vs.\ retraction, live-bet
        authorization) --- Guide 6.
\end{itemize}
\end{notcovered}

\section{What Theseus is}

Theseus is a small research and investment firm whose product is its
\emph{recorded reasoning}, not just its opinions. The three programs
that make that possible are:

\begin{itemize}[leftmargin=*,itemsep=0.4em]
  \item \page{Dialectic} (desktop app). Runs on your laptop during a
        meeting or recording. It listens through the microphone,
        transcribes what people say in near real time, breaks the
        running transcript into \emph{claims} (one-sentence
        assertions), and quietly flags when two claims pull against
        each other. When you stop, it ships an analyzed file straight
        up to the Codex on your behalf.
  \item \page{Noosphere} (the workshop). Takes any recorded material
        --- a Dialectic session, an essay, a PDF, an audio file --- and
        produces an organized view of it: atomic claims, distilled
        principles, a knowledge graph, six-layer coherence checks,
        adversarial objections, methodology profiles, time-tracked
        snapshots. It also runs the live pipelines for Currents and
        Forecasts. You never interact with it directly; it runs against
        material uploaded to the Codex.
  \item \page{Theseus Codex} (website + desktop wrapper). Two faces of
        one site. The \emph{public face} is what readers see at
        \route{theseuscodex.com}. The \emph{founder control plane}
        (signed in) is roughly seventy pages: upload, library,
        knowledge, conclusions, contradictions, explorer, founder
        Currents, founder Forecasts, account, and more.
\end{itemize}

The three pieces share one database, so anything the workshop writes
appears on the Codex immediately, and anything Dialectic uploads gets
handed straight to the workshop.

\section{Two kinds of founder, one shared platform}

It helps to know up front that the platform has \emph{two access
levels}, and most founders will only ever use the first one:

\begin{itemize}[leftmargin=*,itemsep=0.2em]
  \item \textbf{Founder.} Can sign in, upload material, browse and
        review everything the workshop produces, ask the Oracle, watch
        Currents and Forecasts, submit principles for triage, and
        respond to reader engagement. This is what this guide and
        Guides 2--5 cover.
  \item \textbf{Founder-alpha} (operator). All of the above plus the
        operator surfaces: approving conclusions for publication,
        signing the quarterly methodology review week, flipping kill
        switches, authorizing live forecasts, performing deletions or
        retractions, managing roles and credentials. This is Guide
        6's territory. You can read Guide 6 to \emph{understand} what
        founder-alpha does, but the UI for those surfaces will not be
        clickable for you.
\end{itemize}

If you ever land on a page that requires operator privileges, the
Codex shows a clear ``operator only'' notice and links you back to
\page{Dashboard}.

\section{Your forty-minute first session}

\subsection{Step 1 --- sign in to the Codex (3 min)}

Open the Codex URL in your browser. Enter your email and password.
On success you land on \page{Dashboard}. The Codex is two surfaces in
one app: when signed out, visitors see the firm's public face
(reviewed conclusions, articles, Currents, forecasts); once signed in
you see the founder control plane.

The left rail of the founder control plane has many sections. For now
you only need eight, in this order: \page{Dashboard}, \page{Upload},
\page{Library}, \page{Knowledge}, \page{Conclusions}, \page{Ask},
\page{Currents}, \page{Account}.

\screenshot{01_quick_start_dashboard.png}{The authenticated dashboard. The active-principles rail is at the top; recent conclusions and uploads are below; an operations health card sits on the right.}{qs-dashboard}

\subsection{Step 2 --- take a five-minute tour of the dashboard (5 min)}

Read each card on the dashboard without clicking anything yet:

\begin{itemize}[leftmargin=*,itemsep=0.2em]
  \item \textbf{Active principles.} The firm's currently accepted
        principles, conviction-ranked. This is the spine of the firm's
        working memory --- everything the platform does eventually
        cites these.
  \item \textbf{Recent conclusions.} Conclusions the workshop has
        written or updated lately, each with its confidence tier.
  \item \textbf{Recent uploads.} What the firm has been reading.
  \item \textbf{Drift / attention.} Things the system thinks you
        should look at --- a principle whose meaning has been shifting,
        an unresolved contradiction, a source whose standing changed.
  \item \textbf{Operations health.} A quiet indicator that the
        Currents pipeline, embedding index, and forecasts portfolio
        are running. If anything is red, founder-alpha will know
        before you do.
\end{itemize}

\subsection{Step 3 --- install and sign in to Dialectic (5 min)}

Download Dialectic from the firm's releases page and install it.
Launch it. On first run it asks for the Codex URL and your firm
credentials --- the same email and password you used to sign in to the
website. Dialectic does the rest behind the scenes; you do not handle
any API key or token yourself.

Once you've signed in, the credentials are stored locally on your
machine with restrictive file permissions. You only sign in once per
machine.

\subsection{Step 4 --- record a short Dialectic session (10 min)}

Click \page{Begin session}. Talk for five to ten minutes about anything
you have opinions on --- a decision the firm is considering, a thesis
you want to develop, a recent piece of writing. While you talk, the
window shows:

\begin{itemize}[leftmargin=*,itemsep=0.2em]
  \item A running transcript on the left.
  \item A small graph of claims as labeled dots. Each meaningful
        sentence becomes a dot.
  \item A \page{TENSIONS} panel that lights up if two of those claims
        look like they conflict.
  \item A \page{PROMPTS} panel where the optional interlocutor (off
        by default) can suggest a follow-up.
\end{itemize}

When you're done, click \page{Stop session}. Dialectic saves the
analyzed file locally and, because you're signed in, ships it up to
the Codex automatically. The original audio stays on your laptop ---
it is never auto-uploaded.

\screenshot{01_quick_start_dialectic.png}{Dialectic's three-pane window: transcript on the left, claim graph in the middle, tensions and prompts on the right.}{qs-dial}

\subsection{Step 5 --- watch the session land on the Codex (3 min)}

Switch to the Codex. Open \page{Library} (the org-wide upload
inventory) or \page{Upload}. Your new session is there, with a small
\texttt{processing} badge. Click into it. The detail page shows a
live, streaming view of the workshop's progress as it runs: extracting
claims, embedding them, distilling principles, checking coherence,
generating adversarial objections.

When the workshop finishes, the badge clears, the upload's status
becomes \texttt{ingested}, and counts appear next to the row: how many
claims were extracted, how many were grouped into principle
candidates, how many methodology profiles were produced.

\subsection{Step 6 --- upload one document for breadth (3 min)}

Open \page{Upload}. Drag in one written document --- PDF, markdown,
plain text, or a \texttt{.docx}. Hit submit. The file goes through the
same pipeline: extraction, chunking, embedding, claim extraction, etc.
Watch the progress until the row is \texttt{ingested}.

Now you have two sources in the corpus: your spoken session and one
written piece. That's the minimum for everything below to feel real.

\subsection{Step 7 --- explore Knowledge and Conclusions (5 min)}

Open \page{Knowledge}. You will see your uploads listed with their
claim counts. Click into one and browse the extracted claims ---
short, declarative statements, each carrying a pointer back to the
exact span in the source it came from. Click any claim to jump to
that source span.

Open \page{Conclusions}. These are the firm's vetted positions,
including ones the workshop has just drafted from your two uploads.
Each conclusion has a confidence tier (\texttt{firm / founder /
open}), its supporting claims, its source uploads, any contradictions
raised against it, and a \emph{principle shape}: what kind of
principle it is (rule, criterion, mechanism, heuristic, definition,
formula, algorithm), its domain of applicability, and any
quantifiable proxies. Guide 2 takes this surface apart in detail.

\screenshot{01_quick_start_knowledge.png}{The Knowledge surface, showing extracted claims with provenance back to their source spans.}{qs-knw}

\subsection{Step 8 --- ask the Oracle (5 min)}

Open \page{Ask}. Type a question that the material you just uploaded
would plausibly bear on. The Oracle answers in the firm's voice, with
inline citations to the specific claims, conclusions, or principles
it leaned on. Click any citation to jump to the source.

Two things to expect on a first session:

\begin{itemize}[leftmargin=*,itemsep=0.2em]
  \item The Oracle will sometimes \textbf{abstain} (``the corpus does
        not contain enough relevant material to answer this''). That
        is deliberate. The Oracle is built to refuse rather than
        fabricate.
  \item It may show a \textbf{hallucination badge} if its prose
        paraphrases beyond what the citations strictly support. Read
        the cited spans, decide whether the paraphrase is honest,
        narrow your question if needed.
\end{itemize}

Guide 3 explains the full citation contract, the confidence bands,
and how to read these warnings.

\subsection{Step 9 --- glance at Currents (3 min)}

Open \page{Currents}. This is the firm's live commentary surface: if
the Currents pipeline has been running and the corpus has anything
substantive to say about recent X posts, you'll see opinions here,
each grounded in specific firm conclusions. Each opinion's audit trail
shows the source post, which conclusions supported it, the relevance
scores, and the generator metadata. If the corpus is too thin on a
topic, you will see explicit abstentions rather than thin prose.

Currents is also the place to develop a feel for the firm's voice in
the wild. Guide 4 walks the pipeline end to end.

\subsection{Step 10 --- glance at Forecasts (1 min)}

Open \page{Forecasts}. You'll see prediction-market opportunities
(cards) from Polymarket and Kalshi that the firm has scored. Each card
carries a firm probability prediction, a citation chain, a
deterministic decision trace, and a paper or live status. \emph{Live
trading is operator-only}: as a founder you can see what the firm has
forecast and how those forecasts have resolved, but the actual
authorization to place real-money bets goes through founder-alpha.
Guide 5 explains the architecture so you understand what's behind
each card.

\subsection{Step 11 --- where to go next}

You have now touched every primary surface. The pattern to
internalize:

\begin{itemize}[leftmargin=*,itemsep=0.2em]
  \item Material enters through \page{Upload} (manually) or through
        Dialectic's auto-sync (during a conversation).
  \item The workshop processes it, writing \page{Knowledge},
        \page{Conclusions}, methodology profiles, contradictions, and
        principle drafts back into the Codex.
  \item You browse, triage, and respond on the founder control plane.
  \item You ask questions through \page{Ask}, watch \page{Currents}
        for live commentary, and \page{Forecasts} for prediction-market
        action.
  \item Founder-alpha approves what's ready for the world. Once
        approved, the public site picks it up.
\end{itemize}

\section{What you just saw, in one picture}

\begin{center}
\begin{tabular}{l@{\;\(\to\)\;}l@{\;\(\to\)\;}l@{\;\(\to\)\;}l}
\sffamily microphone & \sffamily Dialectic & \sffamily Codex upload row & \sffamily workshop \\
\sffamily upload form & \sffamily Codex upload row & \sffamily workshop & \sffamily knowledge graph \\
\sffamily X / RSS & \sffamily Currents scheduler & \sffamily EventOpinion & \sffamily public feed \\
\sffamily Polymarket & \sffamily Forecasts scheduler & \sffamily Prediction & \sffamily paper or live bet \\
\end{tabular}
\end{center}

Everything that follows in the other five guides refines one of these
four rows.

\section{Common first-session frictions}

\begin{itemize}[leftmargin=*,itemsep=0.3em]
  \item \textbf{Dialectic can't sign in.} Check that the Codex URL is
        the one your firm uses. If your password isn't working, your
        account may need a reset --- ask founder-alpha.
  \item \textbf{Dialectic finishes but the upload doesn't appear in
        the Codex.} Network errors are logged but never block the UI;
        the analyzed file is still on disk. Wait a minute and refresh
        \page{Library}; if nothing shows up, mention it on the firm's
        internal channel.
  \item \textbf{Upload sits in \texttt{processing} for a long time.}
        Long PDFs and audio files take longer. If processing exceeds
        a few minutes, check the dashboard's operations card to see
        whether the workshop is healthy.
  \item \textbf{The Oracle says ``I don't have enough sources.''}
        Correct behavior with a thin corpus. Upload more on the topic,
        or narrow the question. Guide 3 explains what to do with the
        response.
  \item \textbf{No Currents opinions are appearing.} Either the
        ingestion path is paused or the corpus doesn't have enough
        relevant material yet. Either way, founder-alpha will know
        before you do.
\end{itemize}

\section{Where the docs are}

\begin{itemize}[leftmargin=*,itemsep=0.2em]
  \item \file{docs/Theseus\_Codex\_User\_Guide.pdf} --- the older
        single-document Codex walkthrough, still useful as a backup.
  \item \file{dialectic/README.md} --- the Dialectic install notes.
  \item The other five guides in this series.
\end{itemize}

\begin{nextup}
Continue with \href{02_Knowledge_and_Principles.pdf}{Guide 2 ---
Knowledge and Principles} for how the firm turns the material you
just uploaded into reusable principles and vetted conclusions. That
guide picks up the moment your upload's badge clears.
\end{nextup}

\end{document}
